% ============================================================================
%  BRAIn Lab style — компактный референс для Beamer-деки.
%  Палитра: постер (kawasaki/poster/brainposter.sty) + лабовая дека
%  Talk3_Efficient_Pretraining (сливовый фон на титуле и разделах, кремовый
%  на контенте, янтарный текст на сливовом).
%  Build:  latexmk -pdf lab-style-mini.tex
%  Проверка вылезаний:  python ../scripts/check_overflow.py lab-style-mini.pdf
% ============================================================================
\documentclass[aspectratio=169,10pt]{beamer}

\usetheme{metropolis}
\usepackage[T2A]{fontenc}   % [T1] если дека только на английском
\usepackage[utf8]{inputenc}
\usepackage[russian]{babel}
\usepackage{lmodern}
\usepackage{amsmath,amssymb,mathtools}
\usepackage[most]{tcolorbox}
\usepackage{xcolor}
\usepackage{graphicx}
\usepackage{booktabs}

% ---- BRAIn Lab palette -----------------------------------------------------
% Источники: kawasaki/poster/brainposter.sty (постер) + Talk3_Efficient_Pretraining
% (лабовая дека: сливовый фон на титуле/разделах, кремовый на контенте, янтарный текст).
\definecolor{plum}{HTML}{854C65}       % фон титула и разделов; заголовки на креме
\definecolor{cream}{HTML}{FAF8F1}      % фон контентных слайдов; текст на сливовом
\definecolor{amber}{HTML}{FFD183}      % заголовки на сливовом фоне
\definecolor{amberink}{HTML}{A9741A}   % янтарный, читаемый на креме
\definecolor{ink}{HTML}{1F1B18}
\definecolor{red}{HTML}{B0453F}        % проблемы, ограничения
\definecolor{teal}{HTML}{4E7E8C}       % boxcolor постера, затемнённый до читаемого
\definecolor{muted}{HTML}{8C8078}

\setbeamercolor{background canvas}{bg=cream}
\setbeamercolor{normal text}{fg=ink,bg=cream}
\setbeamercolor{frametitle}{fg=plum,bg=cream}
\setbeamercolor{progress bar}{fg=plum,bg=plum!18}
\setbeamercolor{title}{fg=amber}
\setbeamercolor{subtitle}{fg=cream}
\setbeamercolor{itemize item}{fg=plum}
\setbeamercolor{enumerate item}{fg=plum}
\setbeamercolor{alerted text}{fg=plum}
\setbeamercolor{structure}{fg=plum}

\setbeamerfont{title}{series=\bfseries,size=\LARGE}
\setbeamerfont{frametitle}{series=\bfseries,size=\large}
\setbeamertemplate{navigation symbols}{}
\metroset{block=fill,progressbar=none,numbering=none}
% тонкая сливовая линия под заголовком, по ширине текста (не во всю страницу)
\addtobeamertemplate{frametitle}{}{%
  \vspace{-4pt}\par{\color{plum!35}\rule{\textwidth}{0.6pt}}\par\vspace{2pt}}

% полностраничный сливовой фрейм (титул, разделы, финал)
\newenvironment{plumframe}{%
  \setbeamercolor{background canvas}{bg=plum}%
  \setbeamercolor{normal text}{fg=cream}%
  \usebeamercolor[fg]{normal text}%
  \begin{frame}[plain]}{\end{frame}}

% ---- boxes -----------------------------------------------------------------
% fontupper=\small обязателен: иначе три бокса не влезают в колонку.
\tcbset{
  mathbox/.style={fontupper=\small,colback=white,colframe=plum!30,coltext=ink,
    boxrule=0.5pt,arc=2pt,left=5pt,right=5pt,top=3pt,bottom=3pt},
  ideabox/.style={fontupper=\small,colback=plum!7,colframe=plum,coltext=ink,
    boxrule=0pt,leftrule=2.8pt,arc=0pt,left=6pt,right=5pt,top=4pt,bottom=4pt},
  thmbox/.style={fontupper=\small,colback=amber!30,colframe=amberink,coltext=ink,
    boxrule=0pt,leftrule=2.8pt,arc=0pt,left=6pt,right=5pt,top=4pt,bottom=4pt},
  probbox/.style={fontupper=\small,colback=red!6,colframe=red,coltext=ink,
    boxrule=0pt,leftrule=2.8pt,arc=0pt,left=6pt,right=5pt,top=4pt,bottom=4pt},
  figcard/.style={colback=white,colframe=plum!22,boxrule=0.5pt,arc=2pt,
    left=3pt,right=3pt,top=3pt,bottom=3pt},
}
\newcommand{\figcard}[2][\linewidth]{%
  \begin{tcolorbox}[figcard]\centering\includegraphics[width=#1]{#2}\end{tcolorbox}}

\newcommand{\bad}[1]{{\color{red}\bfseries #1}}
\newcommand{\good}[1]{{\color{plum}\bfseries #1}}
\newcommand{\amb}[1]{{\color{amberink}\bfseries #1}}
\newcommand{\ax}[1]{{\color{muted}\scriptsize\rmfamily [arXiv:#1]}}
\newcommand{\fcap}[1]{\par{\color{muted}\scriptsize #1\par}}

% раздел-разделитель: сливовой фон, крупный янтарный заголовок
\newcommand{\topicframe}[2]{%
  \begin{plumframe}
    \vfill
    \begin{center}
      {\color{cream!55}\small #1}\\[0.22cm]
      {\color{amber!50}\rule{0.26\textwidth}{0.7pt}}\\[0.5cm]
      {\color{amber}\fontsize{28}{34}\selectfont\bfseries #2}
    \end{center}
    \vfill
  \end{plumframe}}

\title{Название}
\subtitle{короткий подзаголовок}

\begin{document}

% ===== Титул: только типографика. Логотипа и авторов нет, как в лабовой деке. =
\begin{plumframe}
  \vfill
  \begin{center}
    {\color{amber}\fontsize{40}{46}\selectfont\bfseries Название}\\[0.45cm]
    {\color{amber!50}\rule{0.40\textwidth}{0.7pt}}\\[0.45cm]
    {\color{cream}\large подзаголовок в одну-две строки}
  \end{center}
  \vfill
  {\color{amber!85}\small месяц год}\hfill{\color{amber!85}\small BRAIn Lab}
  \vspace{0.3cm}
\end{plumframe}

\topicframe{01}{Название раздела}

% ===== Разбор работы: фигура из статьи + «что предложили» + «что не так» ======
\begin{frame}{Работа: одна мысль в заголовке}
\begin{columns}[T]
\column{0.45\textwidth}
\figcard{figures/paper-fig.pdf}
\fcap{Что показано на картинке, одна-две строки. \ax{2506.21734}}
\column{0.52\textwidth}
\begin{tcolorbox}[mathbox]
\good{Что предложили.} Формула идёт внутрь бокса, а не отдельно:
\[ z_{i+1}=f_\theta(z_i,x) \]
\end{tcolorbox}
\vspace{0.2cm}
\begin{tcolorbox}[probbox]
\bad{Что не так.} Ограничение работы своими словами, без вежливых оборотов.
\end{tcolorbox}
\end{columns}
\end{frame}

% ===== Широкая фигура сверху + максимум два бокса снизу =====================
\begin{frame}{Когда фигура широкая}
\figcard[0.68\linewidth]{figures/wide-fig.pdf}
\vspace{0.15cm}
\begin{columns}[T]
\column{0.49\textwidth}
\begin{tcolorbox}[mathbox]
\good{Слева.} Под широкой фигурой больше двух боксов не влезает.
\end{tcolorbox}
\column{0.49\textwidth}
\begin{tcolorbox}[ideabox]
\good{Справа.} Вывод, а не пересказ.
\end{tcolorbox}
\end{columns}
\end{frame}

% ===== Факт или теорема (янтарный) + таблица ================================
\begin{frame}{Факт, на который опираемся}
\begin{tcolorbox}[thmbox]
\amb{Утверждение.} Если $\alpha\lambda<1$, то невязка падает линейно:
$\;\|f(z_i)-z_i\|\le(\alpha\lambda)^i\|f(z_0)-z_0\|$.
\end{tcolorbox}
\vspace{0.25cm}
\begin{tcolorbox}[mathbox,fontupper=\scriptsize]
\centering
\begin{tabular}{@{}l@{\hspace{9pt}}l@{\hspace{9pt}}l@{}}
\toprule
{\color{plum}\bfseries метод} & {\color{plum}\bfseries что даёт} & {\color{plum}\bfseries чего не хватает} \\
\midrule
A & короткая ячейка & короткая ячейка \\
B & короткая ячейка & короткая ячейка \\
\bottomrule
\end{tabular}
\end{tcolorbox}
\end{frame}

% ===== Финал ===============================================================
\begin{plumframe}
  \vfill
  \begin{center}
    {\color{amber}\fontsize{34}{40}\selectfont\bfseries Спасибо}\\[0.40cm]
    {\color{amber!50}\rule{0.24\textwidth}{0.7pt}}\\[0.35cm]
    {\color{cream}\large Вопросы?}
  \end{center}
  \vfill
\end{plumframe}

\end{document}
